\section*{Chapitre \ref{ch:s2:the-volatility-index-complex} --- Le complexe des indices de volatilité}

\begin{solution}{exo:s2:the-volatility-index-complex:1}
$A\,L(L - 1)\,r = 100 \times (-1) \times (-2) \times 0.3 = 60$~: il achète des contrats à terme pour une valeur de 60.
\end{solution}

\begin{solution}{exo:s2:the-volatility-index-complex:2}
Avec un indice doublé~: $100 \times 2 \times 1 \times 0.3 = 60$, c'est encore un achat. Avec l'indice seul~: $L(L - 1) = 0$, il ne se passe rien.
\end{solution}

\begin{solution}{exo:s2:the-volatility-index-complex:3}
$10w + 38(1 - w) = 30$ donne $w = 8/28 = 0.286$ sur le premier contrat.
\end{solution}

\begin{solution}{exo:s2:the-volatility-index-complex:4}
La courbe était en contango pendant 89~\% des jours de l'année, si bien que la position longue a payé le roll-down presque chaque jour~; la hausse du VIX est venue par à-coups qui
se sont résorbés, et l'indice a terminé à la hausse alors que les contrats à terme détenus par la position longue avaient tous convergé vers lui à la baisse.
\end{solution}

\begin{solution}{exo:s2:the-volatility-index-complex:5}
Le contrat à terme reflète l'espérance de la racine carrée d'une variance, et la racine carrée est concave~: par l'inégalité de Jensen, l'espérance de la racine est inférieure à la racine
de l'espérance de la variance. L'écart croît avec la volatilité de la volatilité.
\end{solution}

\begin{solution}{exo:s2:the-volatility-index-complex:6}
Leur flux $A\,L(L - 1)\,r$ a le signe du mouvement du jour pour $L = -1$ et $L = 2$~: après une hausse des contrats à terme, ils doivent acheter davantage à la clôture,
ce qui pousse les contrats à terme encore plus haut le jour où la liquidité est la plus mince.
\end{solution}

\begin{solution}{exo:s2:the-volatility-index-complex:7}
L'indice passe l'essentiel de son temps sous sa moyenne (médiane 17.7 contre une moyenne de 19.6), si bien que les contrats à terme, qui anticipent un retour à la moyenne, se traitent au-dessus de lui~: contango pendant 65~\%
des jours, sans aucune prime. Le rendement quotidien moyen de l'indice long est proche de zéro ($-3.7~\%$ par an), mais son rendement logarithmique est de $-26~\%$ par an~:
une série dont la volatilité atteint 88~\% se capitalise très en dessous de sa moyenne.
\end{solution}

\begin{solution}{exo:s2:the-volatility-index-complex:8}
Dans un livre vendeur de volatilité, les rendements ne croissent pas proportionnellement à la taille~: avec la moitié du capital, le livre a été anéanti par le premier krach, et avec le capital entier
plus tôt encore. La hausse de 189~\% de l'indice lors du premier krach dépasse la totalité du capital d'un livre vendeur à hauteur de 0.53 de celui-ci ou plus.
\end{solution}

\begin{solution}{pb:s2:the-volatility-index-complex:1}
\begin{enumerate}
\item L'anticipation du marché pour la volatilité du S\&P 500 à 30 jours, à partir des options SPX et SPXW hors de la monnaie à deux échéances interpolées à 30 jours.
\item Le contrat le plus éloigné comparé au plus proche ou à l'indice~; contango pendant 85.0~\% des jours, le second contrat 0.89 point au-dessus du premier et le premier
de 0.63 au-dessus du VIX.
\item Chaque contrat converge vers l'indice à la baisse à mesure que l'échéance approche~; la position longue a perdu 99.99~\% de 2013 à 2026, soit un taux logarithmique de 66.4~\% par an.
\item 2014~: $-25~\%$ alors que le VIX monte de 35~\%~; 2015~: $-37~\%$ alors que le VIX monte de 2~\%~; 2025~: $-45~\%$ alors que le VIX baisse de 17~\%.
\item Un fonds ou une note qui livre un multiple du rendement quotidien d'un indice de contrats à terme sur le VIX, rééquilibré à la clôture.
\item Valeur $A(1 + Lr)$, contrats à terme détenus $LA(1 + r)$, cible $LA(1 + Lr)$~; l'écart est $A\,L(L - 1)\,r$.
\item Produits inverses ($L = -1$) et produits à effet de levier ($L = 2$).
\item Les produits VIX les plus populaires sont presque assurés de perdre de la valeur au fil du temps~; ceux qui suivent les indices à court terme avaient perdu près de 4~milliards de dollars.
\item VIX de 17.31 à 37.32~; le premier contrat à terme de 15.625 à 33.225.
\item L'indice a progressé de 95.9~\%~; le produit inverse a perdu 95.9~\% le jour même et 96.7~\% depuis le début de l'année.
\item Un événement d'accélération sur le XIV~: sa valeur indicative intrajournalière était tombée à 20~\% ou moins de la clôture précédente.
\item Comme un échec du rééquilibrage de la couverture et de l'effet de levier sur un marché concentré et volatil, à la manière de l'assurance de portefeuille.
\item \textbf{Résultat nommé.} Des contrats à terme vendus à hauteur d'un quart du capital ont rapporté 9.6~\% par an avec un drawdown de 67~\%~; le produit inverse a atteint 16.5 fois
sa valeur et est tombé à zéro lorsque le premier krach a porté l'indice de 17.3 à 63.6, face à un flux égal à 3.77 fois sa valeur à cette clôture.
\item 4.8~\% par an à 0.10 du capital, 9.6~\% à 0.25~; anéanti à 0.50 et à 1.00.
\item Des contrats à terme valant 3.77 fois sa valeur, soit le double des 189~\% du jour.
\item L'indice passe l'essentiel de son temps sous sa moyenne, si bien que le retour attendu fait monter les contrats à terme.
\item La plus forte hausse en un jour que le livre doit survivre, et non son roll-down.
\item Les règlements-livraisons, et non les niveaux de l'indice~; 2018 et 2020 inclus~; le roll et les coûts~; plusieurs tailles.
\item Anticipation des \omterm{def:s2:the-volatility-index-complex:flow}{flux de rebalancement des ETP}.
\item Les acheteurs de contrats à terme sur la volatilité le paient, en échange d'une protection qui rapporte lorsqu'ils en ont besoin.
\end{enumerate}
\end{solution}

\begin{solution}{iq:s2:the-volatility-index-complex:1}
La courbe des contrats à terme est en général en contango~: chaque contrat converge vers l'indice à la baisse à mesure que l'échéance approche, et une position longue à maturité constante ne cesse de
vendre le contrat proche, moins cher, et d'acheter le contrat éloigné, plus cher. Les pics sont brefs~; le roll, lui, est constant.
\end{solution}

\begin{solution}{iq:s2:the-volatility-index-complex:2}
Le marché anticipe une baisse de la volatilité depuis un niveau élevé~: une phase de stress est en cours. Le roll avantage désormais les positions longues~; les positions courtes font face au risque d'un nouveau
pic~; les trades de calendrier et les petites positions longues sont peu coûteux à détenir.
\end{solution}

\begin{solution}{iq:s2:the-volatility-index-complex:3}
Il faut pondérer les deux premiers contrats à chaque clôture par le nombre de jours restants, de sorte que l'échéance soit constante, reporter les pondérations de la veille sur le règlement-livraison du jour, et
valoriser un contrat arrivant à échéance à son règlement final. Les erreurs proviennent des jours d'échéance, des jours fériés, des règlements manquants ou nuls, et de l'utilisation des niveaux de l'indice à la place des
contrats à terme.
\end{solution}

\begin{solution}{iq:s2:the-volatility-index-complex:4}
Il faut qu'il survive à la plus forte hausse en un jour, avec une marge~; tester le stress d'un triplement de l'indice et les flux des produits à la clôture~; limiter la taille
de sorte que la perte en stress ne représente qu'une fraction du capital.
\end{solution}

\begin{solution}{iq:s2:the-volatility-index-complex:5}
À chaque clôture~: la valeur du produit à partir du rendement de l'indice, l'exposition cible $L$ fois la nouvelle valeur, les contrats à terme détenus après le mouvement, et
l'écart à négocier~; avec des contrôles sur les cours de règlement, le calendrier des rolls et un déclencheur d'accélération.
\end{solution}

\begin{solution}{iq:s2:the-volatility-index-complex:6}
Comme établi dans le chapitre~: $LA(1 + Lr) - LA(1 + r) = A\,L(L - 1)\,r$. Sur deux jours, le produit rapporte $(1 + Lr_1)(1 + Lr_2) - 1 = L(r_1 + r_2) +
L^2 r_1 r_2$, alors que $L$ fois le rendement sur deux jours de l'indice vaut $L(r_1 + r_2) + L r_1 r_2$~: l'écart est $L(L - 1)r_1 r_2$.
\end{solution}
